\documentclass[12pt]{article}

\usepackage{url,verbatim, amsmath, amssymb, amsfonts, dsfont, color}
\textwidth 15true cm
\textheight 9true in
\oddsidemargin 0.30true in
\evensidemargin 0.30true in
\topmargin -0.3true in
\headsep 0.4true in

\def\sgn{{\rm sign}}
\def\drightarrow{{\stackrel{d}{\rightarrow}}}
\def\dotsim{\stackrel{.}{\sim}}
\def\goesd{\stackrel{d}{\rightarrow}}
\def\goesp{\stackrel{p}{\rightarrow}}
\def\goeslaw{\stackrel{{\cal L}}{\longrightarrow}}
%
\def\vp{{\varphi}}
\def\E{{\rm E}}
\def\Pvalue{{$p$-value}}
\def\Pvalues{{$p$-values}}
\def\pr{{\rm pr}}
\def\Pr{{\rm pr}}
\def\logit{{\rm logit}}
\def\T{{ \mathrm{\scriptscriptstyle T} }}
\def\diag{{\rm diag}}  
\def\half{{\textstyle{1\over 2}}}
\def\ith{{$i{\rm th}$ }}
\def\var{{\rm var}}
\def\cl{{{\rm c}\ell}}
\def\checkx{{\checked\hspace{-6pt}\setminus}}
\def\grad{{\nabla}}
\def\bs#1{{\boldsymbol{#1}}}

\newcommand{\R}{\mathbb{R}}
\newcommand{\N}{\mathbb N}
\newcommand{\Z}{\mathbb Z}
\newcommand{\Q}{\mathbb Q}
\newcommand{\1}{\mathds 1}
\newcommand{\Var}{\text{Var}}
\newcommand{\Cov}{\text{Cov}}
\newcommand{\sd}{\text{sd}}
\newcommand{\se}{\text{se}}
\renewcommand{\L}{\mathcal {L}}
\renewcommand{\lambda}{\uplambda}
\newcommand{\iidsim}{\stackrel{iid}{\sim}}
\newcommand{\otherwise}{\text{otherwise}}
\newcommand{\indep}{\perp \!\!\! \perp}
\renewcommand{\epsilon}{\varepsilon}
\newcommand{\st}{\text{ s.t. }}
\newcommand{\ds}{\displaystyle}
\newcommand{\ind}{\stackrel{d}{\to}}
\newcommand{\inp}{\stackrel{p}{\to}}
\newcommand{\inas}{\stackrel{a.s.}{\to}}
\DeclareMathOperator*{\argmax}{arg\,max}
\DeclareMathOperator*{\argmin}{arg\,min}

\sloppy %permits line-breaking of urls

\def\blue{\color{blue}}
\def\red{\color{red}}
\def\sgn{{\rm sign}}
\def\drightarrow{{\stackrel{d}{\rightarrow}}}
\begin{document}

\noindent{\bf Questions Week 8} {\hfill STA 2212S 2026}\\

%{\bf Due January 14} \hfill{\it MS, Exercise 5.2}

\bigskip

\begin{enumerate}

\item {\it SM Exercise 9.1.6: } Let $T=0$ with probability $1-\alpha$ and $T=1$ with probability $\alpha, 0<\alpha<1$. Suppose that conditional on $T=0$, $R_0$ follows a normal distribution with expected value $0$ and $R_1$ follows a normal distribution with expected value $\delta$, while conditional on $T=1$, the corresponding expected values are $\eta$ and $\eta + \delta$. In each case the variance of the normal distribution is $1$. Let $$Y = R_0(1-T)+R_1T$$ denote  the  observed response variable. Show that $$\gamma = \E(Y\mid T=1)-\E(Y\mid T=0) = \eta+\delta,$$ and deduce that $\delta = \E(R_1)-\E(R_0)$ cannot be estimated unless $(R_0, R_1)$ and $T$ are independent.

\textcolor{red}{
From the question statement, we have the following conditional distributions:
\begin{align*}
    R_0 \mid T = 0 \sim \mathcal N(0, 1) & \qquad R_0 \mid T = 1 \sim \mathcal N(\eta, 1)\\
    R_1 \mid T = 0 \sim \mathcal N(\delta, 1) & \qquad R_1 \mid T =  1 \sim \mathcal N(\eta + \delta, 1)
\end{align*}
This implies 
\begin{equation*}
    Y \mid T = 1 \sim \mathcal N(\eta + \delta, 1) \qquad Y \mid T = 0 \sim \mathcal N(0, 1)
\end{equation*}
thus 
\begin{equation*}
    \gamma = \eta + \delta - 0 = \eta + \delta
\end{equation*}
Since we only observe $Y$, we can only estimate $\eta + \delta$. However, if $(R_0, R_1)$ are independent from $T$, then $R_i \mid T \stackrel{d}{=} R_i$, so 
\begin{equation*}
    \gamma = E(Y \mid T = 1) - E(Y \mid T = 0) = E(R_1) - E(R_0) = \delta
\end{equation*}
so even though we observe $Y$, we can estimate $\delta$. 
}

\item {\it AoS Exercise 16.6.5: } Suppose that $X \in \mathbb{R}$ is a continuous treatment variable, and that for each subject, the counterfactual regression function is $C_i(x) = \beta_{1i}x$; i.e.~each subject has their own slope. Construct a joint distribution on $(\beta_1, X)$ such that $\Pr(\beta_1 > 0 ) = 1$ but $\E(Y\mid X=x)$ is decreasing in $x$, where $Y=C(X)$. Interpret.

\textcolor{red}{
Let $X \sim \text{Unif}(1, 2)$ and define $\beta_1 = 1/X^2$. Then 
\begin{equation*}
    P(\beta_1 > 0) = P\left(\frac{1}{X^2} > 0\right) = 1
\end{equation*}
since $X^2 \geq 1$. Also, 
\begin{equation*}
    E(Y \mid X = x) = E(\beta_1 X \mid X = x) = E(1/X \mid X = x) = \frac{1}{x} 
\end{equation*}
which is decreasing in $x$. 
}

\end{enumerate}
\end{document}

